\section{Related Work}
\label{sec:related-work}

\noindent\textbf{zkML systems and prover optimization.}
Prior zkML systems reduce proving cost through specialized representations,
protocols, and compilation techniques. ZEN~\cite{boyuan2021zen},
Mystique~\cite{weng2021mystique}, and zkCNN~\cite{liu2021zkcnn} optimize
neural-network proving, while newer systems such as
ZKML~\cite{chen2024zkml}, zkLLM~\cite{sin2024zkllm}, and
ZKPyTorch~\cite{xie2025zkpytorch} extend support to broader workloads.
\TITLE{} is complementary: it leaves the proving protocol unchanged and
focuses on executing the resulting proof work under concurrent requests.

\noindent\textbf{Partitioning and distributed proving.}
DIZK~\cite{wu2018dizk} and Pianist~\cite{liu2024pianist} distribute proof
generation across machines to scale beyond a single prover. Within zkML,
ZKTorch~\cite{chen2025zktorch} exposes parallelism through specialized
basic blocks and parallel proof accumulation, while
DSperse~\cite{ivanov2025dsperse} considers proofs over selected model
slices. Recent orchestration work such as push0~\cite{ahmadvand2026push0}
addresses dispatch and recovery for distributed blockchain proving. These
systems primarily scale an individual proof computation or proving pipeline.
\TITLE{} instead treats proof jobs from concurrent inference requests as a
shared serving workload with scheduling and resource management across the
prover pool.

\noindent\textbf{ML serving systems.}
Clipper~\cite{crankshaw2017clipper}, Clockwork~\cite{gujarati2020serving},
INFaaS~\cite{romero2021infaas}, and Ray Serve~\cite{rayserve} manage
concurrent inference workloads across shared resources. zkML serving differs
in that one request may expand into long-running proof jobs with
stage-varying memory demand and reusable proving artifacts. \TITLE{} applies
serving principles to these proof-specific constraints through request-level
placement and node-local admission.
